\section{Introduction}
\label{sec:introduction}

Pre-training a visual foundation model~\cite{bommasani2021foundation} is
expensive twice over: it takes millions of images and the compute to
match~\cite{caron2021dino,he2022mae}. Self-supervised learning (SSL) removes
the annotation cost by training image representation models on unlabelled
images~\cite{chen2020simclr,zbontar2021barlow}, but the appetite for data and
compute remains, so groups without either --- university labs, small
companies --- can rarely train or even study these models. Making SSL useful
in that setting requires understanding how SSL methods behave when both the
pre-training dataset and the compute budget are small. This work
characterizes one widely used SSL method, Barlow
Twins~\cite{zbontar2021barlow}, in exactly that regime.

We study the question along the data axis. We define the \emph{data
efficiency} of a pre-training method as how its downstream accuracy grows
with the number of unique pre-training images $N$. The pre-training set size
$N$ is our only controlled variable: we deliberately do \emph{not} hold the
number of training steps equal across set sizes, because equal steps would
mean each image in a small set is seen many times more often, turning the
comparison into a compute question rather than a data question. Instead,
every set size trains under a generous fixed epoch budget, and we evaluate
the checkpoint that a standard label-free validation signal selects within
that budget (Section~\ref{sec:method}). We do not claim these models are
trained to convergence; we claim they are trained long enough that the
selected checkpoint is past the point where more training helps, and we
disclose per set size whether that held.

Figure~\ref{fig:pipeline} shows the experimental pipeline. A small vision
transformer (ViT-Tiny with $8\times8$ patches, 5.4M
parameters)~\cite{dosovitskiy2021vit,touvron2021deit} is pre-trained with
Barlow Twins on subsets of Tiny-ImageNet~\cite{le2015tinyimagenet} ranging
from 1\,000 to 100\,000 images. The encoder is then frozen, and a linear
classifier is trained on CIFAR-10~\cite{krizhevsky2009cifar}; its test
accuracy --- the \emph{linear-probe accuracy} --- measures the quality of the
learned representation on a downstream task the encoder never saw. This
pipeline is a deliberately scaled-down version of the foundation-model
recipe: pre-train once on broad data, adapt cheaply to a downstream task.
The study is part of a five-student project in which each student applies a
different SSL method to this shared pipeline (same backbone, same data
subsets, same probe); we chose Barlow Twins, and this paper characterizes
its behaviour.

Barlow Twins is an interesting subject for a small-data study. It learns by
pushing the cross-correlation matrix between the embeddings of two augmented
views of the same images toward the identity matrix, which makes
representational collapse unattractive by construction, without the negative
pairs of contrastive methods~\cite{chen2020simclr} or the asymmetric
momentum/predictor networks of self-distillation
methods~\cite{grill2020byol,caron2021dino}. At the same time, the
joint-embedding family it belongs to has been reported to depend more
strongly on large pre-training sets than masked-prediction
alternatives~\cite{elnouby2021largescale}. Small-data SSL in general is well
studied (Section~\ref{sec:related_work}); Barlow Twins specifically, at this
small a scale and on a small ViT, is not --- its small-data behaviour is both
scarcely described and, given its family's reported data appetite, not
obvious in advance.

We therefore ask: \emph{how does the downstream accuracy of Barlow Twins
scale with pre-training set size for a small vision transformer in a
small-compute regime, and how does it behave at the smallest scales?}

Our contributions are:
\begin{itemize}
  \item A characterization of the Barlow Twins data-efficiency curve
        (Tiny-ImageNet $\rightarrow$ CIFAR-10 linear probe) for a ViT-Tiny
        across $N \in \{1\mathrm{k}, \ldots, 100\mathrm{k}\}$, reported
        against a random-initialization baseline.
  \item An empirical examination of the smallest-data regime: whether Barlow
        Twins trains past its accuracy peak and declines (over-training), and
        whether such a decline coincides with dimensional collapse of the
        representation, measured by effective
        rank~\cite{roy2007effective,jing2022dimensional}.
  \item A study of the interaction between projector width --- the
        method-specific capacity knob of Barlow Twins --- and pre-training
        set size.
\end{itemize}
All findings are observational and scoped to this setup: one backbone, one
pre-training dataset, and one downstream task (Section~\ref{sec:limitations}).

\begin{figure}[t]
  \centering
  % [PLACEHOLDER --- replaced by figures/pipeline.pdf in the figure pass]
  \fbox{\parbox[c][4.5cm][c]{0.92\linewidth}{\centering\small
    Visual abstract (pending): SSL pre-training on Tiny-ImageNet subsets of
    increasing size $\rightarrow$ frozen encoder $\rightarrow$ CIFAR-10
    linear probe $\rightarrow$ one accuracy-vs-$N$ curve per SSL method.}}
  \caption{The data-efficiency pipeline. A ViT-Tiny encoder is pre-trained
  with Barlow Twins on a Tiny-ImageNet subset of size $N$ (no labels), then
  frozen; a linear classifier trained on CIFAR-10 measures the downstream
  accuracy of the learned representation. Repeating this across $N$ from 1k
  to 100k yields a data-efficiency curve: how much downstream accuracy each
  amount of unlabelled pre-training data buys.}
  \label{fig:pipeline}
\end{figure}
